\documentclass[11pt]{article}
\usepackage[T1]{fontenc}
\usepackage{lmodern,amsmath,amssymb,amsthm,booktabs,microtype}
\usepackage[margin=1in]{geometry}
\usepackage[hidelinks]{hyperref}
\newtheorem{theorem}{Theorem}
\newcommand{\OO}{\mathcal O}
\newcommand{\eps}{\varepsilon}
\title{A squarefree composite criterion for a cubic prime race}
\author{AI-written research note by GPT-6.2 Astra}
\date{22 November 2026}
\begin{document}
\maketitle
\begin{abstract}
For the splitting field of $X^3-269$, a one-sided bound for a signed sum over primes tests the critical line of its two-dimensional Artin twists. We replace the primes by squarefree products in $(y,2y]$ of at least two distinct primes, each exceeding a cutoff $z(y)=(\log y)^{2+o(1)}$. The error is $O_{m,\eps}(y^{1/2+\eps})$ unconditionally at each fixed modulus. The sign of the required inequality reverses. We derive the error using periodic cancellation on the Eisenstein lattice and a convolution at an exponent just above $1/2$.
\end{abstract}

\section{The identity and its consequence}
Fix a positive multiple $m$ of three. A prime is called good if it does not divide $3\cdot269m$. For a good prime put
\[
 t(p)=\#\{v\bmod p:v^3=269\bmod p\}-1.
\]
For $p\equiv1\bmod3$, this is $2$ if $269$ is a cube and $-1$ otherwise. For $p\equiv2\bmod3$ it is zero. For $y$ sufficiently large let $z=z(y)\ge2$ satisfy $\log z/\log\log y\to2$ and put
\[
\begin{split}
 P_m(y)&=\sum_{\substack{y<p\le2y\\p\equiv1\bmod m\\p\ \mathrm{good}}}t(p)\log p,\\
 C_m(y)&=\sum_{\substack{y<n\le2y,\ n\equiv1\bmod m\\
                   n\ {\rm squarefree},\ \omega(n)\ge2\\
           p\mid n\,\Rightarrow\,p>z,\ p\ \mathrm{good}}}
                \left(\prod_{p\mid n}t(p)\right)\log n ,
\end{split}
\]
where $\omega(n)$ is the number of distinct prime divisors. All contributing factors of a composite split in the Eisenstein field. The condition modulo $m$ is imposed on the product.

\begin{theorem}\label{th:main}
For every $\eps>0$,
\[
                     P_m(y)+C_m(y)=O_{m,\eps}(y^{1/2+\eps}).
\]
The estimate holds once $y$ is sufficiently large for the chosen cutoff. Fix $\delta>0$. Let $\rho$ be the standard representation of the splitting field of $X^3-269$. The following are equivalent:
\begin{enumerate}
\item All the nontrivial zeros of all $L(s,\rho\otimes\chi)$, $\chi\bmod m$, have real part $1/2$.
\item For every $\eps\in(0,1/2)$ there are $c_{m,\eps,\delta}$ and arbitrarily large $Y$ such that $C_m(y)\le c_{m,\eps,\delta}y^{1/2+\eps}$ for every $y\in[Y,Y^{1+\delta}]$.
\end{enumerate}
\end{theorem}
For $m=3$ this is the RH of $\zeta_{\mathbb Q(\sqrt[3]{269})}/\zeta$. The lower bound for primes and the upper bound for composites carry the same arithmetic difficulty at these exponents.

\section{Cancellation of the full norm coefficients}
Write $\omega=e^{2\pi i/3}$, $\OO=\mathbb Z[\omega]$, $q=269$. The cubic character $\xi$ of $(\OO/q)^\times=\mathbb F_{q^2}^\times$ takes $\alpha$ to the complex lift of $\alpha^{(q^2-1)/3}$. Extend it by zero on multiples of $q$. The character kills all six units because $6\mid(q^2-1)/3$. Since $\OO$ is Euclidean, it defines an ideal character. Let $A(n)=\sum_{N\mathfrak a=n}\xi(\mathfrak a)$. At a prime $p\equiv1\bmod3$, take a primary prime $\pi$ over $p$. Cubic reciprocity gives $(269/\pi)_3=(\pi/269)_3$, so $\xi(\pi)+\xi(\bar\pi)=t(p)$ \cite[Lemma 4]{mer}. At an inert good prime the character is one on the rational generator, and its prime ideal has norm $p^2$. Consequently the good Euler factor of $\sum A(n)n^{-s}$ is $(1-t(p)T+\eta(p)T^2)^{-1}$, for $T=p^{-s}$ and $\eta=\chi_{-3}$. The bad factors have reciprocal polynomials of degree at most two with bounded coefficients.

For every reduced class $b\bmod m$ and $X\ge1$,
\[
 \left|\sum_{\substack{n\le X\\ n\equiv b\bmod m}}A(n)\right|\ll_m\sqrt X . \tag{1}
\]
Indeed set $L=\operatorname{lcm}(m,q)$. Choose a unit $v$ modulo $L\OO$ whose residue modulo $q$ is a noncube, lifting at the $q$-power and using the Chinese remainder theorem at the other primes. The unit $v/\bar v$ has norm one and character $\xi(v)^2\ne1$. Its action permutes every norm fibre modulo $m$ inside $\OO/L\OO$, so the character sum over each fibre vanishes. Each complete period parallelogram in the disk $|\alpha|\le\sqrt X$ has zero sum. The cells meeting the circle cost $O(L\sqrt X+L^2)$ altogether by counting the cells in a boundary strip of width $O(L)$. Dividing by six counts each ideal once. With $m,q$ fixed this proves (1), uniformly in $b$.

\section{Removing repeated and small factors}
For fixed $z$ define $F_z(s)=\prod_{p>z,\ {\rm good}}(1+t(p)p^{-s})$, and write $F_z(s)=\bigl(\sum A(n)n^{-s}\bigr)\sum h_z(d)d^{-s}$ on $\Re s>1$. The local factors of the second series at good primes $p>z$ are shown below. The first-degree coefficient vanishes in every case.
\[
\begin{array}{c|r|r|l}
 \text{permutation of the roots}&t(p)&\eta(p)&(1+t(p)T)(1-t(p)T+\eta(p)T^2)\\ \hline
 \text{identity}&2&1&1-3T^2+2T^3\\
 \text{three-cycle}&-1&1&1-T^3\\
 \text{transposition}&0&-1&1-T^2
\end{array}
\]
At $p\le z$ we use just the reciprocal local polynomial. For $1/2<\sigma<1$, taking absolute values coefficientwise gives
\[
 \sum_{d\ge1}\frac{|h_z(d)|}{d^\sigma}
 \ll_{m,\sigma}
          \exp\!\left(C_\sigma\frac{z^{1-\sigma}}{\log(2z)}\right).
                                                                      \tag{2}
\]
Indeed the large-prime factors are at most $1+3p^{-2\sigma}+2p^{-3\sigma}$ at $T=p^{-\sigma}$, with convergent product bounded independently of $z$. For the small-prime factors use $1+2p^{-\sigma}+p^{-2\sigma}$. Chebyshev's upper bound and partial summation give $\sum_{p\le z}p^{-\sigma}\ll_\sigma z^{1-\sigma}/\log(2z)$. The finitely many bad factors change only the constant.

Write $f_z(n)$ for the coefficients of $F_z$. In $\sum_{n\le X,\ n\equiv1\bmod m}f_z(n)$ every divisor $d$ from $h_z*A$ is coprime to $m$. The complementary factor lies in class $d^{-1}\bmod m$.
When $d\le X$, (1) bounds its partial sum by $O_m((X/d)^{1/2})$, hence by $O_m((X/d)^\sigma)$. The full sum is therefore bounded by $X^\sigma$ times (2). Fix $u=\min(\eps/3,1/4)$ and take $\sigma=1/2+u$. Uniformly in $2\le z\le(\log y)^{2+u}$, the logarithm of the variable factor in (2) is $O_u((\log y)^{1-3u/2-u^2})=o_u(\log y)$. The stated cutoff condition guarantees this range eventually. Throughout $y\le X\le2y$ hold the single value $z(y)$ fixed. Partial summation inserts $\log n$ at cost $O(\log(2y))$. The resulting sum over $(y,2y]$ is $P_m(y)+C_m(y)$ since $z(y)<y$ eventually. Absorbing this logarithm and $y^{o_u(1)}$ into $y^{\eps-u}$ proves the uniform estimate.

The one-sided prime-window theorem \cite{window} gives the last equivalence in Theorem~\ref{th:main}. RH for the indicated twists gives $P_m(y)=O_m(\sqrt y\log^2(2y))$. Conversely, a lower bound $P_m(y)\ge-c_\eps y^{1/2+\eps}$ on the stated windows for each $\eps$ excludes all zeros to the right of $1/2$ by the Landau-window theorem applied to $(2^s-1)(\log\prod_\chi L(s,\rho\otimes\chi))'/s$. A positive real zero gives a pole of positive residue, and the multiplier is nonzero at every nonreal zero in this half-plane. The functional equations reflect the remaining half-strip. Adding the error in Theorem~\ref{th:main} converts the upper composite inequality into precisely that lower prime inequality.

\section*{Finite examples}
The following use $m=3$, $z=(\log y)^2$ and natural logarithms. The first two columns of values on the right count the terms of nonzero weight. The last two give their signed sums, rounded here to two decimals.
\[
\begin{array}{r|rr|rr}
 y&\text{primes}&\text{composites}&P_3(y)&C_3(y)\\ \hline
 12\,000&611&41&306.31&20.04\\
 30\,000&1\,404&144&-178.00&-222.68\\
 75\,000&3\,225&500&-79.96&97.18\\
 150\,000&6\,064&1\,073&-528.86&464.32
\end{array}
\]
For these four finite intervals $z^3>2y$, so the composites can be counted a second way by enumerating unordered pairs of distinct split primes. The file \texttt{check\_rough\_race.py} performs both computations and verifies the three local polynomials. At arbitrarily large $y$, products of three or more primes must also enter the criterion.

\begin{thebibliography}{9}
\bibitem{mer} J.~Merikoski, \emph{A cubic analogue of the Friedlander--Iwaniec spin over primes}, \S2, Lemma 4. \url{https://arxiv.org/abs/2012.05675v2}.
\bibitem{window} \emph{A one-sided prime-interval test for all Dirichlet twists of a cubic Artin factor}, companion note, Theorem 1. Its analytic input is L.~Alp\"oge's \texttt{Deduction.landau\_window} at revision \texttt{7cbf7a0}, in \texttt{landauWindow.lean}, lines 83--84. \url{https://alpo.ge/remark.pdf}.
\end{thebibliography}
\end{document}
